\section{Step 24 Compute signed saturated lane sums}\label{pcm-s24}

Apply \texttt{adds\_s.h} to the original inputs and retain the result in \texttt{w3}. Its signed halfword operation replaces positive and negative overflow with the corresponding endpoints, independently for each lane. \texttt{w2} is still available for the comparison. The register does not carry a permanent halfword type; the instruction’s h suffix selects the lane interpretation.

\Needspace{12\baselineskip}
\subsection{Write and inspect}
Replace the temporary branch and nop in \texttt{vector\_saturate}. Keep the surrounding labels and region markers.

\Needspace{4\baselineskip}
\MARSCodeLabel{MIPS assembly}
\begin{minted}[samepage=true]{mars-idea}
    adds_s.h $w3, $w0, $w1
\end{minted}
Use recipe \texttt{s24-1} from the companion construction inputs. Pause at \texttt{mix\_ready} before comparing this partial state.

\subsection{Expected observation}
The first group’s saturated sums are 0, 3000, -3000, 32767, -32768, 32767, -32768 and 32767.

\Needspace{11\baselineskip}
\subsection{Check your understanding}
Which operation matches signed PCM clipping?

\begin{itemize}
  \item \textbf{A} Unsigned saturation.
  \item \textbf{B} Signed halfword saturation.
  \item \textbf{C} A whole-register bit shift.
\end{itemize}
